\section{Introduction}

Colorectal polyps are abnormal tissue growths that can develop into malignant lesions if they are not detected and removed at an early stage. Colonoscopy is widely used for the examination of the lower gastrointestinal tract and enables clinicians to identify and remove suspicious lesions. Nevertheless, visual inspection remains challenging because polyp appearance varies considerably across patients and imaging conditions. Variations in lesion size, morphology, texture, contrast, and location, together with specular reflections, motion blur, intestinal folds, bubbles, and other visual artifacts, can complicate accurate delineation of polyp regions. Automated polyp segmentation therefore remains an important problem in computer-aided colonoscopy, where reliable pixel-level localization may assist subsequent analysis and clinical decision support.

Accurate polyp segmentation requires a model to simultaneously capture high-level semantic information and preserve sufficiently detailed spatial representations. U-shaped encoder--decoder architectures are well suited to this task because they combine progressively abstract encoder features with higher-resolution representations through skip connections. U-Net established this general design, while subsequent architectures such as UNet++ further demonstrated the importance of multi-level feature interaction and hierarchical decoding \cite{unet,unetpp}. However, repeated downsampling reduces spatial resolution, and fixed decoder operations may not always provide sufficient flexibility for lesions that vary considerably in scale and visual frequency content.

Recent transformer-based and hybrid architectures have improved long-range dependency modeling in medical image segmentation. TransUNet combines convolutional representations with transformer encoding, while Polyp-PVT applies hierarchical transformer features and cross-level aggregation to polyp segmentation \cite{transunet,polyppvt}. These approaches demonstrate the importance of broad contextual modeling, but they can introduce additional architectural complexity. Modern convolutional networks such as ConvNeXt provide an alternative by retaining the efficiency and spatial inductive bias of convolution while incorporating design principles that improve representation capacity. ConvNeXt employs large-kernel depthwise convolution, inverted channel expansion, Layer Normalization, GELU activation, layer scaling, and stochastic depth, making it a strong hierarchical encoder for dense prediction tasks \cite{convnext}.

Multi-scale representation is particularly important in polyp segmentation because lesions may occupy very small regions or extend over a large portion of the image. Multi-Scale Convolution Blocks (MSCBs) address this issue by combining convolutional branches with different receptive fields, allowing a network to aggregate local and broader spatial responses. However, conventional multi-branch aggregation generally combines the branch responses using a fixed formulation that is shared across all input images. Consequently, the relative contribution of different receptive fields is not explicitly conditioned on the frequency characteristics of each individual sample.

Frequency-aware feature processing provides another complementary perspective. Frequency-Aware Feature Enhancement Modules (FAFEMs) decompose deep representations into low- and high-frequency components and use these components to enhance feature responses. Low-frequency representations mainly preserve broader semantic structures, whereas high-frequency representations emphasize local spatial variations and fine structural information. Existing frequency-aware modules are typically used to improve the feature representation at the location where the frequency decomposition is performed. However, the frequency descriptor extracted from deep semantic representations can potentially provide useful global information for guiding later multi-scale processing.

This observation motivates the central idea of this work. Instead of treating bottleneck frequency enhancement and decoder multi-scale refinement as independent operations, we establish an explicit cross-level interaction between them. The frequency descriptor produced by a bottleneck FAFEM is propagated to the first decoder stage and used to regulate the relative contributions of the multi-scale branches in an MSCB. In this way, the receptive-field mixture of the decoder becomes sample dependent rather than fixed across all images.

Directly replacing the original MSCB aggregation with a fully frequency-guided formulation could, however, unnecessarily perturb a stable multi-scale representation. To address this issue, we introduce a bounded residual guidance strategy. The frequency-guided branch aggregation is expressed as a controlled deviation from the conventional MSCB response,

$$
M_{\mathrm{RFG}}
=
M_{\mathrm{ord}}
+
s
\left(
M_{\mathrm{guide}}-M_{\mathrm{ord}}
\right),
$$

where $M_{\mathrm{ord}}$ denotes the ordinary multi-scale aggregation, $M_{\mathrm{guide}}$ denotes the frequency-conditioned aggregation, and $s$ is a bounded learnable guidance-strength parameter. This formulation preserves the conventional MSCB pathway as a reference and allows the model to progressively learn frequency-dependent deviations during optimization.

Based on this mechanism, we develop a ConvNeXt-Tiny U-shaped segmentation framework in which the deepest encoder feature is first enhanced using an existing FAFEM. The resulting low-/high-frequency descriptor is retained and transferred across network levels. After the first decoder feature is fused with the Stage-3 encoder representation, the proposed Residual Frequency-Guided Multi-Scale Convolution Block (RFG-MSCB) uses this descriptor to generate sample-adaptive weights for the multi-kernel MSCB branches. The remaining decoder stages use lightweight conventional fusion and depthwise-separable convolution, keeping the additional architectural modification localized to the first encoder--decoder fusion stage.

The principal contributions of this work are summarized as follows:

\begin{itemize}

\item We introduce a frequency-conditioned multi-scale refinement mechanism, termed RFG-MSCB, which reuses the frequency descriptor generated by an existing bottleneck FAFEM to adaptively weight the multi-kernel branches of an existing MSCB. The contribution lies in the coupling between the frequency descriptor and the multi-scale branches rather than in FAFEM or MSCB themselves.

\item We formulate the frequency-guided aggregation as a bounded residual deviation from the conventional MSCB response. This design preserves the ordinary multi-scale pathway as a reference while allowing a controlled, learnable, sample-dependent modification of branch importance.

\item We introduce cross-level bottleneck-to-decoder frequency guidance by transferring the bottleneck frequency descriptor to the Stage-3 decoder refinement stage. This connects deep semantic frequency information with higher-resolution multi-scale feature processing while leaving the remaining decoder stages lightweight and unchanged.

\item We evaluate the resulting ConvNeXt-Tiny segmentation architecture on multiple public colorectal polyp datasets and perform controlled ablation experiments to examine the contribution of frequency-aware enhancement, residual frequency guidance, and the selected Stage-3 placement.
\end{itemize}
